\section{Prime dilation and nonvanishing}\label{sec:dilation}

The estimates above do not rule out $H_n=0$: the constant term could
cancel the remaining beta terms. We rule this out at indices $n=mp$,
where $m$ is a fixed positive even integer and $p$ is a large prime.
After multiplication by $p^{364}$, the beta coefficients vanish modulo
$p$, while the constant terms reduce to a common nonzero multiple of
their values at $m$. Under the rationality assumption, the remaining
beta terms also vanish modulo all sufficiently large $p$. Hence a
nonzero combined constant at $m$ will imply $H_{mp}\ne0$. We will take
$m=2$.

\begin{samepage}
Write $\Zp$ for the rational numbers whose denominators are prime to $p$.
A \emph{unit} in this ring has both numerator and denominator prime to
$p$, and hence a nonzero reduction modulo $p$. All reductions below are
of quantities proved to be $p$-integral.

\subsection{How a long product reproduces a shorter one}

For example, take $p=5$ and consider $\prod_{j=1}^{10}(5z+j)$.
The factors with $j=5,10$ are
$5(z+1)$ and $5(z+2)$. Removing their two powers of $5$ leaves
\[
 5^{-2}\prod_{j=1}^{10}(5z+j)
 =(z+1)(z+2)\prod_{\substack{1\leq j\leq10\\5\nmid j}}(5z+j).
\]
The remaining factors reduce to the constants $1,2,3,4$, each occurring
twice. Since $1\cdot2\cdot3\cdot4\equiv-1\pmod5$, their product is $1$
modulo $5$. Thus selecting every fifth factor recovers the product of
length two:
\[
 5^{-2}\prod_{j=1}^{10}(5z+j)\equiv(z+1)(z+2)\pmod5.
\]
\end{samepage}
Fractional shifts can change which block is recovered. For example,
in $5z+j+1/3$, divisibility by $5$ means $j\equiv3\pmod5$. Writing
$j=5v+3$ gives
\[
 5z+(5v+3)+\tfrac13=5\bigl(z+v+\tfrac23\bigr).
\]
The shift $1/3$ has become $2/3$; the companion block with shift $2/3$
returns $1/3$. The two blocks exchange, so their product is preserved.

We apply this selection to every numerator and denominator block of
$R_{mp}(pz)$. The factors divisible by $p$ recover the blocks of $R_m$,
and Wilson's theorem evaluates the products of the remaining factors.
Accounting for the powers of $p$ and the normalization gives the
following proposition.

\begin{proposition}\label{prop:dilation}
For fixed positive even $m$ and all sufficiently large primes $p$,
\begin{align}
 p^{364}\rho_i(mp)&\in p^i\Zp \quad(1\leq i\leq364),
                                      \label{eq:rho-dilation}\\
 p^{364}q_b(mp)&\equiv\ch(p)(m+1)^{-365}q_b(m)\pmod p
                        \quad(b=1,2,3).\label{eq:q-dilation}
\end{align}
\end{proposition}

\subsection{The rational-function congruence}

We first prove
\begin{equation}\label{eq:rational-dilation}
 p^{364}R_{mp}(pz)\equiv(m+1)^{-365}R_m(z)\pmod p.
\end{equation}
This is an identity of rational functions over $\Fp$.

For $p>m+1$, both $(mp)!$ and $(mp+1)!$ have valuation $m$.
After division by $p^m$, their units are $(-1)^m m!=m!$ modulo $p$,
by Wilson's theorem and the evenness of $m$. The explicit powers of
$2$, $3$ and $4$ in the normalization reduce to those in $C_m$ by
Fermat's theorem. Thus $C_{mp}$ has valuation $293m$,
and its unit reduces to $C_m/(m+1)^{365}$.

The extra factor $pz-7mp$ in the numerator contributes one factor $p$.
In a block with shift $\theta=a/b$, precisely $18m$ factors are divisible
by $p$. They take the form
\[
 p(z-7m+v+\theta'),\qquad 0\leq v<18m,
 \qquad \theta'=a'/b,\quad 1\leq a'\leq b,\quad pa'\equiv a\pmod b.
\]
The map $\theta\mapsto\theta'$ permutes $\Theta$. When $p\equiv2\pmod3$,
the one-third and two-thirds blocks exchange, leaving their product
unchanged. After division by $p$, the selected factors reproduce
$P_m(z)$. In each block the remaining factors have product
$(-1)^{18m}=1$ modulo $p$.

The divisible denominator factors are $p(z+4y)$, $0\leq y\leq m$,
each with multiplicity $365$. The remaining denominator unit is
$(-1)^{365m}=1$. The net exponent of $p$ is therefore
\[
 293m+(72m+1)-365(m+1)=-364.
\]
The units and remaining factors give exactly \eqref{eq:rational-dilation}.

\subsection{Reducing the partial-fraction coefficients}

We extract the principal parts from \eqref{eq:rational-dilation} by
expanding at each pole. At a selected pole $z=-4y$, put $z=-4y+u$.
The factorization above shows that
\begin{equation}\label{eq:local-integral-series}
 u^{365}p^{364}R_{mp}\bigl(p(u-4y)\bigr)\in\Zp[[u]].
\end{equation}
Multiplication by $u^{365}$ removes the denominator factor at that pole.
After extraction of the powers of $p$ already counted, every other
denominator factor has a unit constant term. Its reciprocal has
$p$-integral power-series coefficients. Reduction therefore commutes
with coefficient extraction, giving in $\Fp[[u]]$
\[
 u^{365}p^{364}R_{mp}\bigl(p(u-4y)\bigr)
 \equiv(m+1)^{-365}u^{365}R_m(u-4y)\pmod p.
\]
The term $a_{i,py}(mp)/(pu)^i$ contributes
$p^{364-i}a_{i,py}(mp)$ to the coefficient of $u^{365-i}$ on the left.
Terms regular at this pole contribute only from degree $365$ onward.
Comparing coefficients yields
\begin{equation}\label{eq:selected-pole}
 p^{364-i}a_{i,py}(mp)\equiv(m+1)^{-365}a_{i,y}(m)\pmod p.
\end{equation}
We must also control the poles not indexed by a multiple of $p$.
Write $h=x+py$ with $1\leq x\leq p-1$ and substitute $t=pz-4x$.
The extra numerator factor is now a unit. Each of the four numerator
blocks has $18m$ divisible factors, while the denominator has $m$
divisible factors, each with multiplicity $365$. The total valuation is
\[
 293m+72m-365m=0.
\]
After removing these powers of $p$, the remaining poles $-4y$ are
distinct modulo $p$.
The same power-series argument, now applied to
$u^{365}R_{mp}(p(u-4y)-4x)$, shows that
\begin{equation}\label{eq:other-pole}
 p^{-i}a_{i,h}(mp)\in\Zp\qquad(p\nmid h).
\end{equation}
Both \eqref{eq:selected-pole} and \eqref{eq:other-pole} imply
$p^{364}a_{i,h}(mp)\in p^i\Zp$; summing over $h$ proves
\eqref{eq:rho-dilation}.

\subsection{Reducing the constant terms}

Apply these coefficient bounds to the formula \eqref{eq:qb} for
$q_b(mp)$. If $p>12m$, every
summation index $j\leq4bmp$ is smaller than $p^2$, so it has valuation
at most one. After multiplication by $p^{364}$, the terms with $p\nmid h$
vanish modulo $p$ by \eqref{eq:other-pole}, even when $p\mid j$.
For $h=py$, terms with $p\nmid j$ vanish by \eqref{eq:selected-pole}.
Only terms with both $p\mid h$ and $p\mid j$ can remain. Write
\[
 h=py,\qquad j=pk,\qquad 1\leq k\leq4by.
\]
The coefficients reduce by \eqref{eq:selected-pole}. For odd $p$,
the weights satisfy $\ch(pk)=\ch(p)\ch(k)$: multiplication by $p$
preserves the two odd residue classes modulo $4$ if $p\equiv1\pmod4$
and exchanges them if $p\equiv3\pmod4$; even $k$ has weight zero.
Taking $\ch(p)$ outside the reduced finite sum gives
\eqref{eq:q-dilation} and completes the proof of
Proposition~\ref{prop:dilation}.
